% ============================================================
% CHAPTER 4 -- SPRINTS 1-2: CORE PLATFORM
% ============================================================
\chapter{Sprints 1--2: Core Platform}
\chaptersommaire
\clearpage

\section*{Introduction}
This chapter presents the first implementation phase of the project. Sprint 1 establishes the
technical base of the platform: authentication, user management, RBAC, backend structure,
database connection, and frontend route protection. Sprint 2 builds the public marketplace
foundation: catalog navigation, asset details, category pages, upload preparation, cart, and
payment flow.

\section{Sprint 1: Authentication, Users and RBAC}

\subsection{Sprint Backlog}

\begin{sprintbox}{Sprint 1 -- Core Access Management}
\textbf{Goal:} implement a secure identity layer and role-based access control so that every
future feature can rely on a stable permission model.
\end{sprintbox}

\begin{table}[H]
\centering
\caption{Sprint 1 Backlog}
\begin{tabularx}{\textwidth}{C{1.1cm} X C{2cm}}
\toprule
\textbf{ID} & \textbf{User Story} & \textbf{Priority} \\
\midrule
US-01 & As a visitor, I can create an account and sign in securely. & High \\
US-02 & As a user, I can access my profile and update my information. & High \\
US-03 & As an administrator, I can assign roles and permissions. & High \\
US-04 & As the system, I can protect API routes and frontend pages according to permissions. & High \\
US-05 & As an administrator, I can trace sensitive access errors in audit logs. & Medium \\
\bottomrule
\end{tabularx}
\end{table}

\subsection{Analysis and Design}

The authentication flow relies on Supabase Auth sessions and JWT access tokens. After login,
the frontend stores the authenticated state and sends the token to the backend through the
\texttt{Authorization} header. The backend middleware verifies the token, loads the user
record, attaches roles and permissions to \texttt{req.user}, and lets protected controllers
continue only if the required permission is present.

The RBAC model is based on six roles: \texttt{member}, \texttt{seller}, \texttt{intern},
\texttt{employee}, \texttt{admin}, and \texttt{super\_admin}. Permissions are stored
separately and linked to roles through \texttt{role\_permissions}, which makes the model
extensible without changing route code.

\subsection{Implementation}

The backend implementation follows the layered structure introduced in Chapter 2:
\texttt{auth.js} defines authentication endpoints, \texttt{authController.js} handles login,
registration, logout, and role requests, while \texttt{userRepository.js} encapsulates SQL
access for user data. Authorization rules are centralized in \texttt{permissions.js} and
then reused by routes such as assets, orders, admin, workspace, chat, and intern requests.

On the frontend, \texttt{AuthContext.jsx}, \texttt{PermissionsContext.jsx},
\texttt{ProtectedRoute.jsx}, and \texttt{Can.jsx} provide route guards and conditional UI
rendering. This avoids duplicating authorization checks in each page.

\subsection{Sprint Review}

At the end of Sprint 1, the platform had a reliable security foundation: users could
authenticate, roles were assigned through the database, frontend pages could be protected,
backend endpoints enforced permissions, and denied access attempts could be traced through
the audit layer.

\section{Sprint 2: Catalog and Marketplace Foundation}

\subsection{Sprint Backlog}

\begin{sprintbox}{Sprint 2 -- Public Marketplace Foundation}
\textbf{Goal:} deliver the first usable public marketplace experience: categories, listings,
asset details, previews, cart, and checkout preparation.
\end{sprintbox}

\begin{table}[H]
\centering
\caption{Sprint 2 Backlog}
\begin{tabularx}{\textwidth}{C{1.1cm} X C{2cm}}
\toprule
\textbf{ID} & \textbf{User Story} & \textbf{Priority} \\
\midrule
US-06 & As a visitor, I can browse all marketplace categories. & High \\
US-07 & As a visitor, I can search and filter assets. & High \\
US-08 & As a buyer, I can consult an asset detail page with previews and metadata. & High \\
US-09 & As a member, I can add assets to a persistent cart. & High \\
US-10 & As a member, I can start checkout and prepare payment. & High \\
US-11 & As a seller, I can access the upload page according to permissions. & Medium \\
\bottomrule
\end{tabularx}
\end{table}

\subsection{Analysis and Design}

The catalog is organized around the \texttt{categories} and \texttt{assets} tables. Each
asset belongs to a category and can have tags, metadata, formats, a thumbnail, a main file,
and status information. The frontend category pages under \texttt{src/Pages/Types} reuse
the same listing logic while applying different category filters.

Asset previews are adapted to the asset type. 3D models use \texttt{ModelViewer.jsx},
textures use \texttt{TextureViewer.jsx} and \texttt{TextureSphere.jsx}, HDRIs use
\texttt{HdriViewer.jsx}, shaders use \texttt{ShaderPreview.jsx}, and scripts use
\texttt{ScriptViewer.jsx}. This separates preview logic from the generic asset detail page.

\subsection{Implementation}

The backend exposes catalog features through \texttt{assets.js}, \texttt{assetController.js},
and \texttt{assetRepository.js}. Queries support listing, filtering, detail loading, download
checks, likes, reviews, and AI summaries. Cart and order preparation are handled through
\texttt{orders.js}, \texttt{orderController.js}, and \texttt{orderRepository.js}.

The frontend implements the public routes in \texttt{App.jsx}: \texttt{/Models},
\texttt{/Textures}, \texttt{/Shaders}, \texttt{/HDRIs}, \texttt{/Audio}, \texttt{/VFX},
\texttt{/Animations}, \texttt{/Environments}, \texttt{/Decals}, \texttt{/Cart},
\texttt{/Checkout}, and \texttt{/View/:id}. Shared services such as \texttt{assetService.js},
\texttt{cartService.js}, and \texttt{orderService.js} isolate API calls from components.

\subsection{Sprint Review}

Sprint 2 delivered the first complete user path: browse assets, open a detail page, preview
the asset, add it to the cart, and start checkout. The work also prepared the later upload,
moderation, and payment features by stabilizing the asset data model.

\section*{Conclusion}
Sprints 1 and 2 delivered the technical and user-facing foundation of IGS Marketplace. The
platform became secure, role-aware, navigable, and ready for the seller and administration
workflows developed in the next phase.
